% Teil II: der durchgerechnete Weg.

\subsection*{Gegeben und gesucht}
Gegeben: $M_E$, $R_E=\qty{6378}{km}$, $G$; für die ISS $h=\qty{400}{km}$;
für geostationär $T$=$24$\,h.
Gesucht: (a) allgemeiner Ausdruck $T(R)$; (b) $v$ und $T$ für die ISS;
(c) $v$ und $R$ für geostationäre Satelliten; (d) Verhalten von $R$ und $T$
bei Abbrennung.

\subsection*{Ansatz}
$a)$ $m v^2/R = GM_Em/R^2$ mit $v=2\pi R/T$ verknüpfen.
$b,c)$ Zahlen einsetzen. $d)$ $E_\mathrm{mech}=-GM_Em/(2R)$; Verlust
$\Rightarrow$ negativeres $E_\mathrm{mech}$ $\Rightarrow$ kleineres $R$,
und aus $T^2\propto R^3$ folgt ein kleineres $T$.

\subsection*{Rechnung}
\textbf{a) Umlaufzeit:}
\begin{equation}
  v = \sqrt{\frac{GM_E}{R}}\qquad\Rightarrow\qquad
  T = \frac{2\pi R}{v} = 2\pi\sqrt{\frac{R^3}{G M_E}}\;.
\end{equation}

\textbf{b) ISS} ($R=R_E+h=6378+400=\qty{6778}{km}$):
\begin{equation}
  v
  = \sqrt{\frac{6.674e-11\cdot 5.974e24}{6.778e6}}\ \si{m/s}
  \approx \qty{7.67e3}{m/s}\approx \qty{7.67}{km/s},
\end{equation}
\begin{equation}
  T
  = \frac{2\pi\cdot 6.778e6}{7.67e3}\ \si{s}
  \approx \qty{5.55e3}{s}
  \approx \qty{92.5}{min}.
\end{equation}

\textbf{c) Geostationär} ($T=24\,\mathrm h=86\,400$\,s):
\begin{equation}
  R
  = \left(\frac{G M_E T^2}{4\pi^2}\right)^{1/3}
  = \left(\frac{6.674e-11\cdot 5.974e24\cdot(8.64e4)^2}{4\pi^2}\right)^{1/3}\ \si{m}
  \approx \qty{4.22e7}{m}
  \approx \qty{42245}{km}.
\end{equation}
Bahnhöhe $=R-R_E\approx\qty{35870}{km}$; Bahngeschwindigkeit
\begin{equation}
  v = \frac{2\pi R}{T}
  = \frac{2\pi\cdot 4.225e7}{8.64e4}\ \si{m/s}
  \approx \qty{3.07e3}{m/s}\approx\qty{3.07}{km/s}.
\end{equation}
Eine geostationäre Position ist nur \emph{über dem Äquator} möglich: nur dort
stimmt die Normale der Bahnebene mit der Erd drehachse überein (sonst hätte
der Satellit eine Nord-Süd-Komponente relativ zur Erdoberfläche).

\textbf{d) Abbrennung durch Restluftreibung.} Die Gesamtenergie
$E_\mathrm{mech}=-GM_Em/(2R)$ nimmt durch Reibungsarbeit ab (negativer). Daraus
folgt, dass $R$ \emph{abnimmt}. Wegen $T^2\propto R^3$ wird die Umlaufzeit
$T$ damit ebenfalls \emph{kürzer}.

\begin{ergebnis}
(a) $T=2\pi\sqrt{R^3/(GM_E)}$.
(b) $v\approx\qty{7.67}{km/s}$, $T\approx\qty{92.5}{min}$.
(c) $R\approx\qty{42245}{km}$ (Höhe $\approx\qty{35870}{km}$), $v\approx
\qty{3.07}{km/s}$; nur über dem Äquator möglich.
(d) Bei Abbrennung nimmt $R$ ab und $T$ wird kürzer.
\end{ergebnis}

\begin{kontrolle}
\textbf{Einheiten:} $T$: $\sqrt{m^3/(kg\cdot s^2\cdot kg)^{-1}\cdot m^3}$ verifiziert
die Sekunde; konkret numerisch $\qty{92}{min}$ für die ISS und $=24$\,h für
geostationäre Bahn als Konsistenzcheck (Annahme, Rechnung).

\textbf{Grenzfall:} sehr kleines $R\to R_E$ ergibt $\approx \qty{84}{min}$
(kürzeste Umlaufzeit), sehr grosses $R$ lange Perioden. Geostationär liegt
$T=24$\,h im erwarteten Bereich.

\textbf{Konsistenz (c):} $T=2\pi\sqrt{R^3/(GM_E)}$ mit $R=4.225e7$\,m gibt
$2\pi\sqrt{(4.225e7)^3/(6.674e-11\cdot5.974e24)}\,\si{s}=8.64e4$\,s $=24$\,h.
\checkmark
\end{kontrolle}

\begin{merksatz}
Kreisbahn: $v=\sqrt{GM/R}$, $T=2\pi\sqrt{R^3/(GM)}$, $E_\mathrm{mech}
=-GMm/(2R)$. Verlust an Energie (Reibung) verkleinert $R$ und verkürzt $T$;
mehr Energieverlust heisst engere, schnellere (aber langsamere lineare
Geschwindigkeit) Bahn.
\end{merksatz}
